\chapter{Alphabet}

\begin{sourcebox}
\centering\itshape
The history of the Cyrillic letters is helpful for learning to read them
\end{sourcebox}

\section{The Cyrillic alphabet}

\begingroup
\small
\setlength{\tabcolsep}{4pt}
\renewcommand{\arraystretch}{1.05}
\begin{longtable}{@{}r >{\cyrillicfont}c p{36mm} p{48mm} p{28mm}@{}}
\toprule
 & \normalfont Printed form & Suggested name (Mathematical Reviews transcription in parentheses) & Approximate sound (but see the next chapter) & Corresponding Greek letter (if any)\\
\midrule
\endfirsthead
\toprule
 & \normalfont Printed form & Suggested name & Approximate sound & Corresponding Greek letter\\
\midrule
\endhead
1 & а & a & as in \emph{ah} & Α alpha\\
2 & б & b & \emph{boo} & \\
3 & в & v & \emph{van} & Β beta\\
4 & г & g & \emph{goo} & Γ gamma\\
5 & д & d & \emph{do} & Δ delta\\
6 & е & e & ye as in \emph{yell} & Ε epsilon\\
7 & ё & --- (ë) & yo as in \emph{York} & \\
8 & ж & zhee (ž) & s in \emph{measure} & \\
9 & з & z & \emph{zoo} & Ζ zeta\\
10 & и & i & \emph{machine} & Η eta\\
11 & й & yot (ĭ) & y in \emph{yet} & \\
12 & к & k & c in \emph{coo} & Κ kappa\\
13 & л & l & \emph{el} & Λ lambda\\
14 & м & m & \emph{em} & Μ mu\\
15 & н & n & \emph{en} & Ν nu\\
16 & о & o & \emph{or} & Ο omicron\\
17 & п & p & \emph{pool} & Π pi\\
18 & р & r & trill & Ρ rho\\
19 & с & s & \emph{ess} & Σ sigma\\
20 & т & t & \emph{too} & Τ tau\\
21 & у & u & oo as in \emph{ooze} & Υ upsilon\\
22 & ф & f & \emph{ef} & Φ phi\\
23 & х & h & ch in \emph{loch} & Χ chi\\
24 & ц & tsee (c) & ts in \emph{tsar} & \\
25 & ч & chee (č) & ch in \emph{cheat} & \\
26 & ш & esh (š) & sh in \emph{shoe} & \\
27 & щ & eshch (šč) & shch in \emph{ash-chap} & \\
28 & ъ & hard sign (") & silent (but see Chap. II) & \\
29 & ы & y & i as in \emph{bill} & \\
30 & ь & soft sign (') & silent (but see Chap. II) & \\
31 & э & reversed e (è) & e as in \emph{ell} & \\
32 & ю & yu (ju) & like the word \emph{you} & \\
33 & я & ya (ja) & like the word \emph{yah} & \\
\bottomrule
\end{longtable}
\endgroup

Only the lower-case letters are given here, since the capitals have the same shape, except for \ru{А} and \ru{Е}, which are as in English, and capital \ru{Б} (No. 2), which is slightly different from lower-case \ru{б}. Italics will be discussed in §10 of Chapter III.

\section{Memorizing the alphabet}

How should a student learn the alphabet? He may be tempted to glance over it once or twice and then proceed to the grammar, on the theory that constant looking at the letters will gradually make him familiar with them. But that method is too slow; it is like trying to be a typist after memorizing the position of the keys. Like the typist, the student of Russian needs mechanical drill, which is here provided in Chapters I and II. The first step is to learn to say the alphabet \emph{a, b, v, \ldots}, simply calling the Cyrillic letters by their English counterparts, as far as possible.

In practicing aloud it is advisable to omit \ru{ё} (No. 7), which in printed Russian is often not distinguished typographically from \ru{е} (No. 6),\corrnote{底本原文为 ``almost never distinguished from e''。现代俄语中省略 \ru{ё} 上两点仍很常见，但 \ru{ё} 是独立字母，在消歧、专名、教学和若干规范性场合会明确写出；故仅将 ``almost never'' 改为 ``often not''。} to say \emph{yot} for \ru{й} (No. 11), which is often called ``i-short'' but which we shall always regard as a consonant, namely the English consonantal \emph{y} (but see II§9); and to say \emph{hard sign} and \emph{soft sign} for No. 28 and No. 30, \emph{reversed e} for \ru{э} (No. 31), and \emph{you, yah} for \ru{ю, я} (Nos. 32, 33). For letters like \ru{ж} = ž (No. 8) etc., with the diacritical mark in the Mathematical Reviews transcription, one may either say \emph{zee-check} etc. or else invent names of one's own like \emph{zhee} for \ru{ж} (No. 8), \emph{esh} for \ru{ш} (No. 26), \emph{tsee} for \ru{ц} (No. 24), imitating the sound.

In order to learn the alphabet quickly one should repeat it on all possible occasions, at home and elsewhere, aloud and silently, rapidly and slowly, noting that it falls naturally into two groups: \emph{a, b, v, g, d, e} and \emph{i, yot, k, l, m, n, o, p, r, s, t, u} (note the absence of \emph{q}), which are separated by \emph{zhee, z} (instead of \emph{f, g, h} as in English) and are followed by twelve letters at the end, in four sets: (\emph{f, h}), (\emph{c, č, š, šč}), (hard sign, \emph{y}, soft sign), and (reversed $e$, $yu$, $ya$).

The letters should also be written frequently, with any convenient fixed order for the strokes, both with pen or pencil on paper and with the index finger of the right hand on the palm of the left; and they should be visualized as printed on the wall, on the ceiling and elsewhere.

\begin{keybox}[Exercise 1.1]
The names of these fifteen Russian mathematicians are to be pronounced aloud and copied out in Cyrillic, each letter being named as it is written. The list is arranged in order of date of birth, and a field of special interest is given for each mathematician.
\end{keybox}

\begin{center}
\small
\begin{tabularx}{0.86\textwidth}{@{}>{\cyrillicfont}X r l@{}}
\toprule
\normalfont Name & Birth & Field\\
\midrule
Лузин & 1883 & set theory\\
Хинчин & 1894 & number theory\\
Урысон & 1898 & topology\\
Голузин & 1906 & complex variables\\
Натансон & 1906 & real variables\\
Фукс & 1907 & complex variables\\
Курош & 1908 & algebra\\
Михлин & 1908 & functional analysis\\
Наймарк & 1909 & functional analysis\\
Канторович & 1912 & functional analysis\\
Линник & 1915 & probability\\
Рохлин & 1919 & topology\\
Вишик & 1921 & differential equations\\
Яглом & 1921 & differential geometry\\
Дынкин & 1924 & probability\\
\bottomrule
\end{tabularx}
\end{center}



\section{History of the Cyrillic consonants}

Since the history of the consonants is less important for understanding the present-day system of Russian writing than the history of the vowels (see next section), the following information is given in smaller type. The reader should extract from it whatever help he can in his practical task of learning the alphabet.

The first question naturally is: why does the Cyrillic alphabet begin with letters pronounced like A, B, V and written \ru{А, Б, В}? The answer is that in Greek the pronunciation of the letter Β (beta), which in classical Greek represented a sound like English B, changed gradually during the post-classical and Byzantine periods toward the sound /v/.\corrnote{底本把这一历时变化精确写成 ``about 200 A.D.'' 的单一时点；实际是逐步演变，故只改正时间性质。} The B-sound consequently ceased to be represented by Β in later Greek, while in the Slavic languages this sound was common, so that in Cyrillic a separate letter \ru{Б} (No.~2) was required for it. Greek itself could represent /b/ in later borrowings by the sequence μπ; this provides an instructive comparison, but the shape of Cyrillic \ru{Б} is not securely established as a simple combination of the Greek letters Μ and Π.\corrnote{底本把 \ru{Б} 直接解释为 Μ 与 Π 的组合，并以现代希腊 μπ 为依据；该字形来源不能这样确定，故仅修正这一断言。} The component strokes are more clearly visible in older Cyrillic forms than in the civil type introduced under Peter the Great in 1708.

We now come to a question of considerable general interest: why does the Cyrillic \ru{Н} (No.~15), with the same sound as an American N, look like an ``aitch''? We shall see that the present shape of the vowel \ru{И} = i (No.~10) is thereby explained as well.

The Greek vowel Η (eta) was pronounced in medieval Greek like the i in \emph{machine} and was therefore taken over in Cyrillic as the ancestor of \ru{И} (No.~10). In medieval Cyrillic hands its form developed into the familiar diagonal-stroked \ru{И}. The consonant \ru{Н}, originally taken over from Greek Ν to represent the same sound as English n, developed its modern crossbar form in the independent graphic history of Cyrillic. The two developments can be compared, but they should not be understood literally as one crossbar ``sliding'' in one direction while the other letter moved in the opposite direction to fill a gap.\corrnote{底本用两条横画相向 ``slip'' 的故事解释 \ru{И/Н} 字形。这个说法过于机械；正文保留其比较思路，同时把字形史改为较稳妥的表述。}

In Italy, on the other hand, Latin H continued, through Etruscan and western Greek alphabetic traditions, a sign for /h/; its history is not explained by saying that Greek Η became ``available'' merely because Latin writing did not distinguish long and short vowels.\corrnote{底本将 Latin H 的来源解释为 Greek eta 因“不再需要表示长元音”而转作 /h/；这混淆了不同希腊字母传统及伊特鲁里亚中介，故直接改正。} Since the sound of English h is rather different from the Greek Χ and therefore from Russian \ru{Х} (No.~23), traditional Russian spellings of a number of established Western names used \ru{Г} (No.~4) for historical h, as in older renderings of \emph{Hamlet}; present-day transcription practice is not uniform and often uses \ru{Х}.\corrnote{底本概括为 ``modern Russian usually transliterates an English H with Г''。这更适用于若干传统译名而非现代一般规则，故限定其范围。}

The differences in shape of other consonants appearing in both Latin and Cyrillic are relatively easy to understand; e.g. the Cyrillic \ru{С} (No.~19) and the Latin S are both historically related to Greek sigma. But five consonant letters were necessary in Cyrillic to represent sounds not occurring in Greek; namely \ru{Ж} (No.~8), \ru{Ц} (No.~24), \ru{Ч} (No.~25), \ru{Ш} (No.~26), and \ru{Щ} (No.~27).

The graphic history of these five letters is more complicated than a set of simple direct borrowings. Their forms are related to the development of the early Slavic alphabets, including Glagolitic and Cyrillic, and some show ultimate or comparative connections with non-Greek graphic traditions; but it is not established that \ru{Ж} was simply taken from Coptic, that \ru{Ц, Ш} and probably \ru{Ч} were directly copied from Hebrew, or that \ru{Щ} is merely a ligature of two \ru{Ш}'s.\corrnote{底本把这些 Coptic/Hebrew/ligature 来源逐一写成确定事实。现代文字史对此更谨慎；故对确定性断言作最小事实性修正。}

The sound represented by \ru{Ш} (No.~26), for example, nevertheless has an entertaining linguistic history. A similar /š/-sound is found in ancient Egyptian and in Phoenician and Hebrew, as well as in English and the Slavic languages such as Bulgarian and Russian, but it was not a phoneme for which classical Greek or Latin required an ordinary alphabetic letter (which is why English writes it with two letters \emph{sh}). When the Greek alphabet was constructed from Phoenician, signs unnecessary for Greek phonology were repurposed or discarded; many centuries later the Slavic alphabets again required a distinct sign for /š/. Thus the comparison still helps to explain why Cyrillic contains letters for Slavic sounds absent from Greek, even though the individual graphic derivations are not as simple as the older account suggested.

\section{History of the vowel-symbols; the basic vowel-scheme}

But it is the Cyrillic system for writing vowels that raises the most interesting and instructive questions. There are ten vowel-letters, forming the basic vowel-scheme.

\begin{center}
\begin{tabular}{ccccc}
\ru{у} & \ru{о} & \ru{а} & \ru{э} & \ru{ы}\\[2pt]
\ru{ю} & \ru{ё} & \ru{я} & \ru{е} & \ru{и}
\end{tabular}
\end{center}

First we must note that although this scheme contains ten different letters, the Russian language has in the traditional elementary analysis used here only five distinct vowel-sounds (more accurately, six);\corrnote{底本把这一计数直接当作唯一分析。现代标准俄语常分析为五个元音音位，并对 [ɨ]/\ru{ы} 与 /i/ 的关系有不同分析；这里保留底本的教学计数但明确其性质。} for except that \ru{ы} and \ru{и} in the last column are somewhat different from each other, the two vowel-symbols in any given column, namely \ru{у} and \ru{ю}, \ru{о} and \ru{ё}, \ru{а} and \ru{я}, \ru{э} and \ru{е}, are intended in this scheme to have the same basic vowel sound. The difference between them is chiefly orthographic: after a consonant the letters in the second row \ru{ю, ё, я, е, и} normally signal its softness, while \ru{ю, ё, я, е} also represent /j/ plus a vowel word-initially, after a vowel, and after separating \ru{ъ/ь}.\corrnote{底本说第二行字母一概“written after soft consonant sounds”。这忽略了词首、元音后及分隔符号后的 /j/+元音功能，故只修正这一正字法事实。} Thus the six vowel qualities distinguished for present purposes may be described as: the u-sound (Column One), the o-sound (Column Two), the a-sound (Column Three), the e-sound (Column Four), and the two i-sounds (Column Five).

The five columns in the vowel-scheme are arranged in the order ``back to front''; i.e. the a-sound is the central vowel, for which the tongue remains comparatively flat (the reason why a doctor asks his patient to say ``ah''); u is the ``back vowel'', for which the tongue is chiefly in the back of the mouth; and for the other vowels the tongue moves generally forward to an extreme frontal position for the sound of i in \emph{machine}. (Here again the reader is urged to experiment aloud.) Moreover, a rough distinction between ``back'' and ``front'' articulations applies also to consonants, as can be realized by experimenting with back consonants such as k and g and more front consonants such as p, b, t, d.

Since softening a consonant consists of bringing the tongue toward the hard palate (see next section), it is more natural for front vowels (columns 4 and 5) to occur after soft consonants than after hard consonants. This helps explain the distribution which the vowel-scheme is intended to summarize. As might be expected, several of the common vowel-letters were taken over directly from Greek alphabetic forms: \ru{у} (No.~21) from upsilon, \ru{о} (No.~16) from omicron, \ru{а} (No.~1) from alpha, \ru{е} (No.~6) from epsilon, and \ru{и} (No.~10; historically connected with Greek eta). The history of the remaining five letters, however, is more complicated than the simple constructions given in the older account: \ru{Ю} is an early Slavic iotated vowel letter; modern \ru{Я} developed through the history of earlier Cyrillic nasal/iotated letters and later civil type; \ru{Э} was standardized in the Russian civil alphabet of the early eighteenth century; \ru{Ё} emerged in late-eighteenth-century Russian orthographic practice; and \ru{Ы} historically arose from a combination involving a yer and an i-letter.\corrnote{底本把 \ru{Ю} 解释为 iota+omicron+alpha，把 \ru{Я} 解释为 Greek alpha 的直接改形，把 \ru{Ё} 的两点简单定于1797年，并把 \ru{Ы} 描述为 iota+hard sign；这些字母史说明过于简化或不准确，故仅在这一段直接修正。}

The letter \ru{Ъ} (No.~28), now called hard sign, and the soft sign \ru{Ь} (No.~30) historically represented the two reduced vowels known as yers. When weak yers disappeared from spoken East Slavic, their alphabetic functions changed. The Russian spelling reform of 1917--1918 abolished the routine final \ru{ъ} after hard consonants and removed several obsolete letters; it did not abolish \ru{ъ} and \ru{ь} themselves, which remain where needed for separating, palatalization and morphological functions (see Chapter II).\corrnote{底本说这些声音消失后，两字母也于1918年 ``were dropped from the writing''，仿佛只作为例外残留；实际改革主要取消词末硬音符并废除若干其他字母，\ru{ъ/ь} 本身仍属俄文字母，故直接改正。}

\section{Hard and soft consonants}

Just as English has more than 26 different sounds to be represented by the 26 letters of our Latin alphabet, so in the elementary count used in this book Russian is described as having 42 sounds (six vowels and 36 consonants) to be represented by the 33 Cyrillic letters.\corrnote{“42 = 6+36”是底本采用的教学计数；现代音系分析的音位数取决于分析框架，故加限定而不改本书体系。} However, Russian spelling, unlike English, is comparatively regular: the pronunciation of a Russian word can usually be inferred from its spelling together with stress and a small number of phonetic rules. There are nevertheless real exceptions, including unstressed-vowel reduction, voicing assimilation and consonant clusters containing unpronounced segments.\corrnote{底本写 Russian spelling ``is almost phonetic''，并称 ``With negligible exceptions, there are no silent letters''；这一说法过强，故只改正这一句。} In certain situations (see below) the consonant \ru{й} (No.~11) is pronounced though not written in the elementary description used here.

The consonant sounds of Russian fall for the purposes of this book into two classes: hard and soft, and are represented in the Cyrillic alphabet by 21 letters. Three of these 21 consonants, \ru{ц} (No.~24), \ru{ж} (No.~8), \ru{ш} (No.~26), are normally hard, and three of them, \ru{й} (No.~11), \ru{ч} (No.~25), \ru{щ} (No.~27), are normally soft. The consonants represented by the other 15 letters may be either hard or soft.

The distinction between hard and soft consonants can be approximately illustrated in English. The words \emph{moo} and \emph{mew} differ in the palatal gesture following or accompanying the initial consonant, but the English sequence /m+j/ is not identical with the single Russian palatalized consonant $/m^{j}/$.\corrnote{底本直接称 \emph{mew} 的 m 为“soft m”；英语通常分析为辅音+/j/序列，不能与俄语 $/m^{j}/$ 完全等同，故把例子明确为近似。} As we have seen, one purpose of the letters in the second row of the vowel-scheme is to indicate softness of a preceding consonant. Thus in \ru{тук} (tuk) nitrate the \ru{т} is hard (as in \emph{toot} in English), but in \ru{тюк} (tyuk) package it is soft; the two Russian words have the same basic /u/ vowel. Soft consonants will be studied in detail in the next chapter.

Moreover, these second-row vowel-symbols (except \ru{и}; see just below) are used after consonants to indicate softness, but they also occur word-initially, after vowels, and after separating \ru{ъ/ь}, where they represent an initial /j/-sound before the vowel.\corrnote{底本说它们 ``are written only after a soft-consonant sound''，若无书面辅音则“理解出一个未写的 \ru{й}”；这不是完整的现代正字法描述，故直接改正为两类环境。} Thus the Russian word \ru{а} but is pronounced like the English \emph{ah}, but the Russian word \ru{я} myself is pronounced with an initial /j/, approximately like English \emph{yah}; and similarly for \ru{ю, ё, е}. But \ru{и} (No.~10) should be pronounced ee, not yee, i.e. the Russian word \ru{и} and is pronounced ee. Note that softness and hardness refer only to consonants and are nevertheless indicated orthographically in part by the choice of vowel-symbol, \ru{у} or \ru{ю}, \ru{а} or \ru{я}, etc. This characteristic feature of Cyrillic writing reflects the historical development of Slavic palatalization and iotation.

\section{The spelling rule}

Since softening of a consonant consists in bringing the tongue toward the hard palate, it is natural, as we have seen, for soft consonants to occur before front vowels and hard consonants before back and central vowels, and in earlier Russian this distribution was much more restrictive than it is today. In modern Russian regular spelling restrictions still hold especially after the three back consonants \ru{к, г, х} (Nos.~12, 4, 23) and after the sibilants; combinations which violate the native spelling rules occur chiefly in foreign words.

After the sibilant consonants \ru{ц, ж, ш} (normally hard) and \ru{ч, щ} (normally soft), where it is unnecessary to indicate hardness or softness in the same way, the spelling of the vowels is restricted. In particular, after \ru{ж, ш, ч, щ}, the u-sound, a-sound and i-sound are written \ru{у, а, и}, not \ru{ю, я, ы}; \ru{ц} has its own distribution, with \ru{ы} occurring in some endings and lexical items.\corrnote{底本先把 \ru{ц} 与其余四个咝音一起讨论，随后规则却把 \ru{ц} 排除，并只笼统说 \ru{ы} ``usually'' after \ru{ц}。这里保留原教学规则，同时把 \ru{ц} 的独立分布说清楚。} Thus the word for fluid is written \ru{жидкость} but pronounced with a hard \ru{ж}; in similar positions \ru{ю} appears as \ru{у} and \ru{я} as \ru{а}.

So we have the ``spelling-rule'':
\begin{sourcebox}
The \emph{u}-sound, the \emph{a}-sound or the \emph{i}-sound, occurring in a native Russian word after a back consonant \ru{к, г, х} or a sibilant \ru{ж, ш, ч, щ}, is never spelled with \ru{ю, я} or \ru{ы} but always with \ru{у, а} or \ru{и}.
\end{sourcebox}

This rule is important in the inflection of nouns, adjectives and verbs (see Chapter III).

\begin{keybox}[Exercise 1.2]
The following names of non-Russian mathematicians born before the nineteenth century are to be pronounced aloud, in essentially the same way as in their original language except that the \ru{г} used for \emph{h} (e.g. \ru{Гопф} for Hopf) is pronounced as \emph{g}. Note also the \ru{эй} for German \emph{eu} in \ru{Эйлер} for Euler.\footnotemark
\end{keybox}
\footnotetext{\small\textbf{校勘：}底本生年有多处误植，现直接校正并保持按出生年排序：Girard Desargues 1591（原1593）；Abraham de Moivre 1667（原1661）；Jacopo Riccati 1676（原1616）；Leonhard Euler 1707（原1701）；Pierre-Simon Laplace 1749（原1736）；Carl Friedrich Gauss 1777（原1771）；Marc-Antoine Parseval 1755（原1805，并从练习1.3移入本表）。Isaac Newton 按英格兰当时旧历为1642年12月25日、公历为1643年1月4日，故记1642/43。}

\begingroup
\footnotesize
\setlength{\tabcolsep}{3.5pt}
\renewcommand{\arraystretch}{1.04}
\begin{longtable}{@{}>{\cyrillicfont}p{29mm} p{21mm} r @{
\hspace{7mm}} >{\cyrillicfont}p{29mm} p{21mm} r@{}}
\toprule
\normalfont Name & Field & Birth & \normalfont Name & Field & Birth\\
\midrule
\endfirsthead
\toprule
\normalfont Name & Field & Birth & \normalfont Name & Field & Birth\\
\midrule
\endhead
Фибоначчи & no. theor & 1170 & Эйлер & anal & 1707\\
Кардано & alg & 1501 & Клеро & diff equ & 1713\\
Виет & alg & 1540 & Лагранж & calc var & 1736\\
Галилей & math phys & 1564 & Лаплас & math phys & 1749\\
Мерсенн & no. theor & 1588 & Маскерони & geom & 1750\\
Дезарг & geom & 1591 & Лежандр & ser & 1752\\
Декарт & geom & 1596 & Менье & diff geom & 1754\\
Кавальери & calc & 1598 & Парсеваль & ser & 1755\\
Уоллис & calc & 1616 & Фурье & ser & 1768\\
Паскаль & geom & 1623 & Гаусс & anal & 1777\\
Ньютон & math phys & 1642/43 & Пуассон & diff equ & 1781\\
Лейбниц & calc & 1646 & Бессель & math phys & 1784\\
Бернулли & calc & 1654 & Брианшон & geom & 1785\\
Лопиталь & calc & 1661 & Навье & math phys & 1785\\
Муавр (де) & compl var & 1667 & Понселе & geom & 1788\\
Риккати & diff equ & 1676 & Коши & diff equ & 1789\\
Тейлор & ser & 1685 & Кориолис & math phys & 1792\\
Стирлинг & ser & 1692 & Грин & math phys & 1793\\
Маклорен & calc & 1698 & Ламе & math phys & 1795\\
Крамер & alg & 1704 & Штейнер & geom & 1796\\
 & & & Штаудт & geom & 1798\\
\bottomrule
\end{longtable}
\endgroup

\begin{keybox}[Exercise 1.3]
Similarly for these non-Russian mathematicians born in the nineteenth century. Note the \ru{а} for \emph{u} in \ru{Рассел} for Russell.\footnotemark
\end{keybox}
\footnotetext{\small\textbf{校勘：}Charles Hermite 生于1822（非1854）；本表已直接改正生年并移至相应年代位置。}

\begingroup
\footnotesize
\setlength{\tabcolsep}{3.5pt}
\renewcommand{\arraystretch}{1.04}
\begin{longtable}{@{}>{\cyrillicfont}p{29mm} p{21mm} r @{
\hspace{7mm}} >{\cyrillicfont}p{29mm} p{21mm} r@{}}
\toprule
\normalfont Name & Field & Birth & \normalfont Name & Field & Birth\\
\midrule
\endfirsthead
\toprule
\normalfont Name & Field & Birth & \normalfont Name & Field & Birth\\
\midrule
\endhead
Абель & alg & 1802 & Фуко & math phys & 1819\\
Штурм & diff equ & 1803 & Кейли & alg geom & 1821\\
Якоби & anal & 1804 & Эрмит & no. theor & 1822\\
Гамильтон & calc of var & 1805 & Кронекер & alg & 1823\\
Дирихле & no. theor & 1805 & Кельвин & math phys & 1824\\
Грассман & geom & 1809 & Риман & diff geom & 1826\\
Лиувилль & diff equ & 1809 & Кристоффель & diff geom & 1829\\
Куммер & no. theor & 1810 & Дедекинд & no. theor & 1831\\
Сильвестр & no. theor & 1814 & Липшиц & diff equ & 1832\\
Вейерштрасс & anal & 1815 & Бельтрами & diff geom & 1835\\
Стокс & math phys & 1819 & Жордан & alg & 1838\\
Мах & math phys & 1838 & Картан & diff geom & 1869\\
Гиббс & math phys & 1839 & Борель & set theor & 1871\\
Ли & groups & 1842 & Рассел & logic & 1872\\
Шварц & compl var & 1843 & Каратеодори & compl var & 1873\\
Кантор & set theor & 1845 & Леви-Чивита & diff geom & 1873\\
Миттаг-Леффлер & compl var & 1846 & Лебег & real var & 1875\\
Клейн & geom & 1849 & Ландау & no. theor & 1877\\
Фробениус & groups & 1849 & Харди & no. theor & 1877\\
Хевисайд & math phys & 1850 & Ритц & math phys & 1878\\
Линдеман & alg & 1852 & Фубини & real var & 1879\\
Лоренц & math phys & 1853 & Эйнштейн & math phys & 1879\\
Риччи-Курбастро & diff geom & 1853 & Рис & func anal & 1880\\
Пуанкаре & diff equ & 1854 & Эддингтон & math phys & 1882\\
Стилтьес & calc & 1856 & Данжуа & real var & 1884\\
Пеано & anal & 1858 & Лефшец & topol & 1884\\
Чезаро & ser & 1859 & Литлвуд & no. theor & 1885\\
Уайтхед & logic & 1861 & Банах & func anal & 1892\\
Гильберт & func anal & 1862 & Гопф & topol & 1894\\
Адамар & anal & 1865 & & & \\
\bottomrule
\end{longtable}
\endgroup
